% Verbatim copy of /Users/alingalatan/workspaces/grothe/handoff/manuscript/legacy/ch25-parts/89b-mixed-contact.tex (file dated 2026-10-25). Only this line was added.
% New dimension-free original-response branch, 22 October 2026.
\subsection{Mixed cubic concentration: two integrals and four endpoints}
\label{sec:mixed-contact-concentration}

\paragraph{What changes.}
The preceding diagonal-Wick cone misses even a pure mixed monomial.
The following \st{P,C} theorem controls that direction with an arbitrary
coupled remainder. It concerns actual bounded responses, not raw K.
It assumes neither a fresh buffer nor a moment inequality, and does not
change the complete Grothendieck endpoint.

Keep $\epsilon=1891/2500$, $p=2500/1891$, $Ph=tZ$, $F=\epsilon Qk$,
$V=\|F\|^2$, and $\mathcal N$ from
\S\ref{sec:contact-concentration-cone}. For standard complex Gaussian
linear forms $X,Y$ orthogonal to each other and to $Z$, define
\[
 J_{X,Y}=\frac{X^2\bar Y}{\sqrt2},\qquad
 c_{X,Y}=|\langle J_{X,Y},F\rangle|.
\]
The three test coordinates are independent. The remainder of $F$ is
\emph{not} assumed independent of any of them.

\begin{theorem}[\st{P,C}: mixed concentration alternative]
\label{thm:mixed-contact-concentration}
In every finite dimension, throughout $21/25\le t\le851/1000$,
\[
 \boxed{\begin{aligned}
 c_{X,Y}^2\ge\frac9{10}V
   &\quad\Longrightarrow\quad \mathcal N<\frac{14259}{20000}=0.71295,\\
 c_{X,Y}^2\ge\frac{179}{200}V
   &\quad\Longrightarrow\quad \mathcal N<\frac{28529}{40000}=0.713225.
 \end{aligned}}
\]
Both conclusions are below the design cap $35663/50000<500/701$.
The second hypothesis is equivalently
$F=\zeta J_{X,Y}+R$, $R\perp J_{X,Y}$, and
$\|R\|^2\le(21/179)|\zeta|^2$.
There is no amplitude restriction on $\zeta$; $R$ may use $X,Y,Z$ and
arbitrarily many other coordinates, with all visible sectors and phases.
\end{theorem}

\subsubsection{The same support lemma, with a mixed test}
Multiply $J_{X,Y}$ by a unit phase so that
$u=\langle X^2\bar Y,F\rangle=\sqrt2c_{X,Y}\ge0$.
Such a phase change preserves the test law. For
\[
 L=Z\{\ell+b(|Z|^2-2)\},\quad D=dX^2\bar Y,\quad
 I(\ell,b,d)=\E|L+D|,
\]
the two signs have equal expectation by $Y\mapsto-Y$.
The pointwise two-disk support inequality and $D\in\operatorname{ran}Q$
give, with $m=\Re\langle Z(|Z|^2-2),h\rangle$,
\[
 \ell t+bm+pdu\le I(\ell,b,d).
\]
Consequently, if $I\le H$ and $b>0$,
\begin{equation}\label{eq:mixed-contact-penalty}
 \mathcal J:=\|Pk\|^2+\epsilon\|Qh\|^2
 \ge\frac{[\ell t+pdu-H]_+^2}{2pb^2}.
\end{equation}
Indeed $\|Z(|Z|^2-2)\|^2=2$, so
$\epsilon\|Qh\|^2\ge m^2/(2p)$.
For $b=0$ the support inequality is a feasibility cut.
The unknown remainder never enters these pairings.
As before, this is one lower bound for the existing retained penalty;
it must not be added to another bound spending the same square.

Only two support integrals are needed:
\[
\begin{array}{c|c|c|c}
 \ell&b&d&I(\ell,b,d)<H\\\hline
 1&0&1/4&467607/500000\\
 1&37/200&87/500&33463/40000
\end{array}
\]
Fresh directed upper values on the default mesh are below
$0.935213488681$ and $0.836574516962$, respectively.
The rational bounds, not these rounded explanatory decimals, are used.

\subsubsection{A closed formula for the concentration-dependent cap}
Put
\[
 t_-=21/25,\quad t_+=851/1000,\quad u_0=47/200,\quad
 H_0=467607/500000,\quad d_0=1/4,
\]
and let $(b,d,H)=(37/200,87/500,33463/40000)$. Define
\[
 U(t)=\frac{H_0-t}{pd_0},\qquad
 A(t,u)=t+pdu-H,
\]
\[
 C_\eta(t,u)=t^2+\frac{pu^2}{2\eta}
                   -\frac{A(t,u)^2}{2pb^2},\qquad
 L_\eta(t)=t^2+\frac{pu_0^2}{2\eta}
                   -\frac{((19/10)t-I_0)^2}{2p},
 \quad I_0=124233393/10^8.
\]
For every $0<\eta\le1$, the same proof supplies the canonical envelope
\begin{equation}\label{eq:mixed-contact-canonical-cap}
 \boxed{\quad c_{X,Y}^2\ge\eta V\quad\Longrightarrow\quad
 \mathcal N\le\max\{L_\eta(t_+),C_\eta(t_+,u_0),
                  C_\eta(t_-,U(t_-)),C_\eta(t_+,U(t_+))\}.\quad}
\end{equation}
Each entry is an explicit rational affine function of $1/\eta$.
Thus the concentration threshold need not be found by another
parameter-grid computation.

Here is the complete reduction. Concentration and the first support cut
give $V\le u^2/(2\eta)$ and $u\le U(t)$.
For $u\le u_0$, use the visible support test
$Z(|Z|^2-1/10)$, whose expectation is less than $I_0$.
This bounds $\mathcal N$ by $L_\eta(t)$, increasing on the central interval.
Its integral is recomputed, rather than imported, from
\[
 \E\sqrt T\,|T-1/10|
 =\frac75\frac{\sqrt\pi}{2}
   +2(1/10)^{5/2}\sum_{n\ge0}
       \frac{(-1/10)^n}{n!(n+3/2)(n+5/2)}<I_0.
\]
For $u\ge u_0$, $A(t,u)>0$ and the second trial gives
$\mathcal N\le C_\eta(t,u)$.
Its coefficient of $u^2$ is
$p/(2\eta)-pd^2/(2b^2)>0$, so the endpoints $u_0,U(t)$ suffice.
The first endpoint increases in $t$ on the interval. The second is
convex in $t$, with quadratic coefficient
\[
 1+\frac1{2\eta p d_0^2}
     -\frac{(1-d/d_0)^2}{2pb^2}>0.
\]
All these signs hold for $0<\eta\le1$; the two involving $1/\eta$
are weakest at $\eta=1$ and are checked there with fractions.
The remaining signs are independent of $\eta$.
This proves \eqref{eq:mixed-contact-canonical-cap}.

The four endpoints, in the displayed order, are below
\[
\begin{array}{c|rrrr}
 \eta=9/10&.711700923&.712935712&.712879259&.712935618\\
 \eta=179/200&.711927522&.713162312&.713219783&.713202005
\end{array}
\]
and prove Theorem~\ref{thm:mixed-contact-concentration} by exact rational
comparison. The checker also checks the strict comparisons with $500/701$.

\subsubsection{Product compression: a derivative-free integral certificate}
Let $M(a,r)=\E_\theta|a+re^{i\theta}|$.
In the three independent Rayleigh radii $x=|Z|$, $y=|X|$, $z=|Y|$,
\[
 I(\ell,b,d)=\E M(A(x),d y^2z),\qquad
 A(x)=(\ell-2b)x+bx^3.
\]
Here $\ell\ge2b$ for both trials. The integrand is convex separately
in $A$, $y^2$ and $z$. On each radial cell in $[0,6]$, replace each
amplitude distribution by its two endpoint atoms of the same mass and
first moment. This is convex domination, so it requires no derivative
bound near zeros of the norm.

The visible and radius moments use the directed primitives from the
previous subsection. The square-amplitude moment is especially simple:
\[
 \int_x^y2r^3e^{-r^2}\,dr
       =(x^2+1)e^{-x^2}-(y^2+1)e^{-y^2}.
\]
The resulting discrete product amplitudes are compressed a second time:
replace every $r=y^2z$ by its exact barycentric mixture between adjacent
nodes on a rational amplitude grid. Convexity of $M(A,dr)$ makes this
another upper bound, preserving mass and first moment. Compression
reduces the three-dimensional tensor sum to a two-dimensional sum.
It does not approximate the unknown remainder's distribution.

For the angular integral use the positive series
\[
 M(a,r)=m\sum_{n\ge0}a_n w^n,\quad
 m=\max(a,r),\quad w=(\min(a,r)/m)^2,
\]
\[
 a_0=1,\quad a_n=a_{n-1}\frac{(2n-3)^2}{4n^2},\quad
 \sum a_n=4/\pi.
\]
Replace the tail after index $N$ by
$(4/\pi-\sum_{n=0}^Na_n)w^{N+1}$; this is an upper enclosure.
Set the value at $a=r=0$ to zero.
All primitives, Machin's formula for $\pi$, roots and sums are enclosed
with integer fixed-point intervals of scale $10^{40}$.

For completeness, with $a=\ell-2b$, $E=e^{-36}$,
$J_2=(73/12)E$, $J_4=(1801/8)E$, the full discarded tail is at most
\[
 aJ_2+bJ_4+E\{38d+2(a+3b/2)\}+dJ_2.
\]
This follows by summing the three radial-tail events and using
$M(A,d y^2z)\le A+d y^2z$,
$\E A\le a+3b/2$, $\E y^2=1$, $\E z\le1$,
and $\E[y^2\mathbf1_{y>6}]=37e^{-36}$.
The tails in the two trials are below $1.3\cdot10^{-14}$.

\paragraph{Survivors and the remaining obstruction.}
Together with the previous Wick cone, any central response exceeding
the design cap must satisfy both
\[
 \sup_{X\perp Z}|\langle H_X,F\rangle|^2<\frac45V,
 \qquad
 \sup_{X,Y\perp Z,\,X\perp Y}|\langle J_{X,Y},F\rangle|^2
                  <\frac{179}{200}V.
\]
The new branch is not contained in the diagonal-Wick branch:
for the unit tensor $J_{X,Y}$ the largest diagonal-Wick share is only
$4/27$. Neither condition implies small influence. For example the
unit tensor $X_1X_2\bar X_3$ has maximal diagonal-Wick share $2/27$
and maximal mixed-square share $1/2$.
Indeed for $X=\sum a_jX_j$, $Y=\sum b_jX_j$ the latter pairing has
squared modulus $2|a_1a_2\bar b_3|^2\le1/2$; equality uses
$X=(X_1+X_2)/\sqrt2$, $Y=X_3$.
The diagonal pairing has squared modulus
$2|a_1a_2a_3|^2\le2/27$, also attainable.
Thus there is still no exhaustive concentration/diffusion alternative,
and the flank rows at the same $\epsilon$ still need coverage.

\paragraph{Replay and formalization.}
The package \code{grothendieck-mixed-contact-20261022.zip} contains a
standard-library \code{verify\_all.py}, the proof and Markdown handoff,
fresh default/fine logs, exact endpoint values and a SHA256 manifest.
Default: radial mesh $1/64$, 32 angular terms; fine: mesh $1/128$,
48 terms, and twice the product-grid resolution.
The Lean port requires the support lemma, the normalized mixed cubic
projection, convex endpoint-atom domination, the two integral bounds,
and four rational endpoint comparisons. It requires no new hypothesis
on the actual cubic's law. This region must be combined by taking the
minimum of valid caps, not by adding retained-penalty savings.
